\section{Related Work}

\paragraph{LSM compaction and tuning.}
The original LSM-tree design buffers updates and later merges them into
disk-resident sorted components~\cite{oneil1996lsm}. The survey by Luo and
Carey describes the resulting design space~\cite{luo2020lsmsurvey}, and
Sarkar et al.\ formalize compaction triggers, layout, granularity, and data
movement as design primitives~\cite{sarkar2021lsmspace}. Monkey and
Dostoevsky optimize Bloom-filter memory, merge policy, and space-time
tradeoffs~\cite{dayan2017monkey,dayan2018dostoevsky}; WiscKey and PebblesDB
change value placement or merge structure~\cite{lu2016wisckey,raju2017pebblesdb}.
bLSM, TRIAD, and Accordion improve write smoothing, logging, and memory
organization~\cite{sears2012blsm,balmau2017triad,bortnikov2018accordion}.
SILK schedules internal LSM I/O to reduce latency spikes, while ADOC, vLSM, and
ArceKV tune dataflow, LSM shape, or compaction thresholds for a single engine
or workload~\cite{balmau2019silk,yu2023adoc,xanthakis2024vlsm,liu2026arcekv}.
Dremel adaptively tunes \rocksdb parameters with an online bandit on a single
instance~\cite{zhao2022dremel}, and Dong et al.'s \rocksdb retrospective at
Meta tracks how amplification priorities shifted with hardware and workload,
still inside one engine~\cite{dong2021rocksdbevolution}. SAKI is orthogonal:
it leaves each engine's compaction policy and parameters unchanged, treats
these single-engine contributions as the inner control loop, and asks how a
fixed host budget should be divided across co-located instances as demand
drifts. The two control layers can be combined.

\paragraph{Multi-tenant storage control.}
Argon, PARDA, mClock, Pisces, Cake, and IOFlow allocate or route storage
service across shared servers, hypervisors, and software-defined storage
planes~\cite{wachs2007argon,gulati2009parda,gulati2010mclock,shue2012pisces,
wang2012cake,thereska2013ioflow}. Linux cgroups and IOCost provide lower-layer
block-I/O weighting~\cite{linux_cgroup_v2,heo2022iocost}, and Gimbal extends
the line to SmartNIC JBOFs that schedule co-tenant NVMe traffic on the
storage data path~\cite{min2021gimbal}. These systems enforce fairness or
SLOs over I/O streams but do not observe LSM debt signals or choose which
\rocksdb instance receives background compaction service. SAKI uses those
engine-level signals and treats kernel or network-side I/O control as a
companion enforcement layer rather than the placement mechanism itself.

\paragraph{Multi-tenant LSM deployments.}
Recent KV-store work moves isolation closer to the storage engine. FlexEngine
uses admission control for multi-tenant serverless LSM deployments, Iso-KVSSD
and ZenFS+ use device namespace or zoned-placement mechanisms for isolation,
and Delta Fair Sharing proposes in-engine multi-tenant fairness for
RocksDB~\cite{wang2026flexengine,min2021isokvssd,oh2023zenfs,
griggs2026deltafairsharing}. CaaS-LSM and DFlush move compaction or flush work
into disaggregated or DPU-assisted settings~\cite{yu2024caaslsm,ding2025dflush}.
SAKI instead targets the common deployment where tenants are separate,
unmodified RocksDB instances on one host and the available background budget is
fixed.

\paragraph{RocksDB controls.}
RocksDB exposes per-instance options, statistics, event listeners, and rate
limiting~\cite{rocksdb,rocksdb_ratelimiter}. The auto-tuned rate limiter solves
local adaptation within a ceiling; SAKI supplies the host-level rank-and-assign
loop that places budgets across instances under a conserved aggregate budget.
Public CacheLib and Baleen traces inform our trace-qualification audits, but
cache admission, prefetching, or request-trace replay are separate from the
fixed-budget RocksDB compaction-placement problem~\cite{cachelib_trace,
baleen_fast24}. SAKI reallocates a conserved host budget across separate
instances without changing the engine.
